% Chapter 5: General Requirements
\chapter{一般要求对比}

\section{安装基本要求}

\subsection{总体原则对比}

\begin{table}[H]
\centering
\caption{安装基本要求对比}
\small
\begin{tabularx}{\textwidth}{|l|X|X|X|}
\hline
\textbf{项目} & \textbf{\versionAShort} & \textbf{\versionBShort} & \textbf{\versionCShort} \\
\hline
安装完整性 & 车辆装备齐全 & 强化，新增"对自动驾驶的安装完整性要求" & 强调完整性与功能性 \\
\hline
安装稳定性 & 应牢固安装 & 强化抗振动要求 & Installation stability \\
\hline
防松措施 & 一般要求 & 新增"扭矩规定"具体数值参考 & 防松措施具体化 \\
\hline
互换性 & 同型号可互换 & 强化 & 可互换性 \\
\hline
热管理 & 一般要求 & 新增"对 LED 灯具的热管理" & 热管理专门化 \\
\hline
\end{tabularx}
\end{table}

\subsection{具体条款对比}

\subsubsection{安装位置（基于实际法规文本）}

\begin{table}[H]
\centering
\caption{安装尺寸要求对比（基于实际法规文本，单位：mm）}
\small
\begin{tabularx}{\textwidth}{|X|X|X|X|}
\hline
\textbf{项目} & \textbf{\versionAShort (2019)} & \textbf{\versionBShort (草案 V4.3.5)} & \textbf{备注} \\
\hline
近光灯横向 & 离外缘端面距离 $\le 400$，两灯间距 $\ge 600$（车宽 $<1300$ 时 $\ge 400$） & 同左 & M1/N1 类无最小间距要求 \\
\hline
近光灯高度 & 500--1200（N3G 可至 1500）& 同左 & 一致 \\
\hline
前雾灯横向 & $\le 400$ & 同左 & 一致 \\
\hline
前雾灯高度 & $\ge 250$；M1/N1 $\le 800$，其他 $\le 1200$，N3G $\le 1500$ & 同左 & 一致 \\
\hline
角灯高度 & 250--900 & 同左 & 一致 \\
\hline
角灯纵向 & 离车前 $\le 1000$ & 同左 & 一致 \\
\hline
转向灯高度（1/1a/1b/2a/2b 类） & $\ge 350$，$\le 1500$（放宽至 2100） & 同左 & 一致 \\
\hline
转向灯高度（5/6 类） & M1/N1 $\ge 350$，其他 $\ge 500$，$\le 1500$（放宽至 2300） & 同左 & 一致 \\
\hline
侧转向灯距离 & $\le 1800$（M1/N1 放宽至 2500） & 同左 & 一致 \\
\hline
制动灯 S1/S2 高度 & 350--1500（放宽至 2100） & 同左 & 一致 \\
\hline
制动灯 S3/S4 位置 & 高于 S1/S2；$\ge 150$（相对后窗下边缘）或 $\ge 850$ 离地 & 同左 & 一致 \\
\hline
前位灯高度 & 250--1500（O1/O2 放宽至 2100） & 同左 & 一致 \\
\hline
后位灯高度 & 350--1500（放宽至 2100） & 同左 & 一致 \\
\hline
后雾灯高度 & 250--1000（N3G 或组合时 $\le 1200$） & 250--1000（与车尾灯组合时 $\le 1200$，N2G/N3G/M2G/M3G $\le 1400$） & 草案略有放宽 \\
\hline
驻车灯高度 & M1/N1 无；其他 350--1500（放宽至 2100） & 同左 & 一致 \\
\hline
示廓灯横向 & $\le 400$（放宽至 500） & 同左 & 一致 \\
\hline
非三角形后回复反射器高度 & 250--900（组合时 $\le 1200$） & 同左 & 一致 \\
\hline
三角形后回复反射器高度 & 250--900（组合时 $\le 1200$） & 同左 & 一致 \\
\hline
非三角形前回复反射器高度 & 250--900（不能保证时 $\le 1500$） & 同左 & 一致 \\
\hline
侧回复反射器高度 & 250--900（组合时 $\le 1200$） & 同左 & 一致 \\
\hline
侧标志灯高度 & 250--1500（不能保证时 $\le 2100$） & 同左 & 一致 \\
\hline
昼间行驶灯高度 & 250--1500 & 同左 & 一致 \\
\hline
倒车灯高度 & 250--1200 & 250--1200（N2G/N3G/M2G/M3G $\le 1400$） & 草案略有放宽 \\
\hline
牵引杆挂车标志灯高度 & 离前挡板 300--400 & 同左 & 一致 \\
\hline
慢行灯高度 & 无规定 & 侧面 $\le 1500$；后部 不高于车辆 & 草案新增 \\
\hline
ADS 标志灯高度 & 未规定 & $\ge 250$ & 草案新增 \\
\hline
\end{tabularx}
\end{table}

\subsubsection{几何可见度（基于实际法规文本）}

\begin{table}[H]
\centering
\caption{几何可见度对比}
\begin{tabular}{|l|c|c|c|}
\hline
\textbf{装置类型} & \textbf{\versionAShort} & \textbf{\versionBShort} & \textbf{\versionCShort} \\
\hline
近光灯可见度      & 横向 $\pm 45°$、纵向 $\pm 15°$ & 同左 & 横向 $\pm 45°$、纵向 $\pm 15°$ \\
\hline
远光灯可见度      & 横向 $\pm 30°$（部分车型 $\pm 20°$） & 同左 & 横向 $\pm 30°$ \\
\hline
DRL 可见度        & 横向 $\pm 20°$ & 横向 $\pm 25°$（强化） & 横向 $\pm 20°$ \\
\hline
转向灯可见度      & 横向 $\pm 40°$（部分 $\pm 80°$） & 同左 & 横向 $\pm 40°$ \\
\hline
\end{tabular}
\end{table}

\section{颜色与色度要求}

\subsection{各装置颜色规定}

\begin{table}[H]
\centering
\caption{各装置颜色规定对比}
\small
\begin{tabularx}{\textwidth}{|l|X|X|X|}
\hline
\textbf{装置类型} & \textbf{\versionAShort} & \textbf{\versionBShort} & \textbf{\versionCShort} \\
\hline
远光灯          & 白色           & 同左，新增"色度容差详细化" & White \\
\hline
近光灯          & 白色                    & 同左 & White \\
\hline
前雾灯          & 白色或淡黄色             & 同左，新增色度边界详细规定 & White or selective yellow \\
\hline
后雾灯          & 红色                    & 同左 & Red \\
\hline
倒车灯          & 白色                    & 同左 & White \\
\hline
前转向灯        & 琥珀色                  & 强化色度容差 & Amber \\
\hline
后转向灯        & 琥珀色                  & 同左 & Amber \\
\hline
侧转向灯        & 琥珀色                  & 同左 & Amber \\
\hline
危险警告        & 琥珀色                  & 同左 & Amber \\
\hline
制动灯          & 红色                    & 同左 & Red \\
\hline
ESS              & 红色（高强度闪烁）      & 强化规定，明确闪烁模式 & Red (high frequency) \\
\hline
前位置灯        & 白色                    & 同左 & White \\
\hline
后位置灯        & 红色                    & 同左 & Red \\
\hline
牌照灯          & 白色                    & 同左 & White \\
\hline
DRL              & 白色                    & 同左，新增色度范围详细化 & White \\
\hline
侧标志灯        & 琥珀色（前）/ 红色（后） & 同左 & Amber/Red \\
\hline
前示廓灯        & 白色                    & 同左 & White \\
\hline
后示廓灯        & 红色                    & 同左 & Red \\
\hline
后反射器        & 红色                    & 同左 & Red \\
\hline
侧反射器        & 琥珀色                  & 同左 & Amber \\
\hline
前反射器        & 白色                    & 同左 & White \\
\hline
\end{tabularx}
\end{table}

\subsection{色度坐标对比（基于实际法规文本）}

\textbf{白色光色度坐标（CIE 1931）—— 装置发出的光色（GB 4785-2019 表 1）}

\begin{table}[H]
\centering
\caption{白色色度区域边界（CIE 1931，源自 GB 4785-2019 表 1）}
\begin{tabular}{|l|l|}
\hline
\textbf{边界编号} & \textbf{边界方程} \\
\hline
W12 绿色边界  & $y = 0.150 + 0.640x$ \\
\hline
W23 黄绿色边界 & $y = 0.440$ \\
\hline
W34 黄色边界  & $x = 0.500$ \\
\hline
W45 紫红色边界 & $y = 0.382$ \\
\hline
W56 紫色边界  & $y = 0.050 + 0.750x$ \\
\hline
W61 蓝色边界  & $x = 0.310$ \\
\hline
\end{tabular}
\end{table}

\begin{table}[H]
\centering
\caption{白色色度区域边界交点（CIE 1931）}
\begin{tabular}{|c|c|c|}
\hline
\textbf{交点} & \textbf{$x$} & \textbf{$y$} \\
\hline
W1 & 0.310 & 0.348 \\
\hline
W2 & 0.453 & 0.440 \\
\hline
W3 & 0.500 & 0.440 \\
\hline
W4 & 0.500 & 0.382 \\
\hline
W5 & 0.443 & 0.382 \\
\hline
W6 & 0.310 & 0.283 \\
\hline
\end{tabular}
\end{table}

\textbf{选择性黄色光色度坐标（CIE 1931）}

\begin{table}[H]
\centering
\caption{选择性黄色色度区域边界（CIE 1931）}
\begin{tabular}{|l|l|}
\hline
\textbf{边界编号} & \textbf{边界方程} \\
\hline
SY12 绿色边界  & $y = 1.290x - 0.100$ \\
\hline
SY23 光谱轨迹   & --- \\
\hline
SY34 红色边界  & $y = 0.138 + 0.580x$ \\
\hline
SY45 黄白色边界 & $y = 0.440$ \\
\hline
SY51 白色边界  & $y = 0.940 - x$ \\
\hline
\end{tabular}
\end{table}

\textbf{琥珀色光色度坐标（CIE 1931）}

\begin{table}[H]
\centering
\caption{琥珀色色度区域边界（CIE 1931）}
\begin{tabular}{|l|l|}
\hline
\textbf{边界编号} & \textbf{边界方程} \\
\hline
A12 绿色边界  & $y = x - 0.120$ \\
\hline
A23 光谱轨迹   & --- \\
\hline
A34 红色边界  & $y = 0.390$ \\
\hline
A41 白色边界  & $y = 0.790 - 0.670x$ \\
\hline
\end{tabular}
\end{table}

\begin{table}[H]
\centering
\caption{琥珀色色度区域边界交点（CIE 1931）}
\begin{tabular}{|c|c|c|}
\hline
\textbf{交点} & \textbf{$x$} & \textbf{$y$} \\
\hline
A1 & 0.545 & 0.425 \\
\hline
A2 & 0.560 & 0.440 \\
\hline
A3 & 0.609 & 0.390 \\
\hline
A4 & 0.597 & 0.390 \\
\hline
\end{tabular}
\end{table}

\textbf{红色光色度坐标（CIE 1931）}

\begin{table}[H]
\centering
\caption{红色色度区域边界（CIE 1931）}
\begin{tabular}{|l|l|}
\hline
\textbf{边界编号} & \textbf{边界方程} \\
\hline
R12 黄色边界  & $y = 0.335$ \\
\hline
R23 光谱轨迹   & --- \\
\hline
R34 紫线       & --- \\
\hline
R41 紫色边界  & $y = 0.980 - x$ \\
\hline
\end{tabular}
\end{table}

\begin{table}[H]
\centering
\caption{红色色度区域边界交点（CIE 1931）}
\begin{tabular}{|c|c|c|}
\hline
\textbf{交点} & \textbf{$x$} & \textbf{$y$} \\
\hline
R1 & 0.645 & 0.335 \\
\hline
R2 & 0.665 & 0.335 \\
\hline
R3 & 0.735 & 0.265 \\
\hline
R4 & 0.721 & 0.259 \\
\hline
\end{tabular}
\end{table}

\textbf{说明：} 三版本在装置发出的光色色度坐标上保持一致。GB 4785 草案 V4.3.5 和 UN R48-09 沿用相同的色度区域定义。

\textbf{反射器光色（GB 4785-2019 表 2 与表 3）}

表 2（回复反射后的夜间光色）和表 3（装置反射的日间光色）规定了反射器的特殊色度要求，与光源色度有所不同。

\textbf{关键差异点：}

\begin{enumerate}[leftmargin=2em]
    \item 白色光色度边界在三版本中保持一致，使用相同的 6 条边界方程和 6 个交点；
    \item 琥珀色边界方程在 GB 4785-2019 中使用 $y = x - 0.120$（绿色边界），UN R48 中也采用相同定义；
    \item 红色边界在 GB 4785-2019 中下沿为 $y = 0.335$，与 UN R48 保持一致；
    \item 反射器色度（夜间与日间）有独立的色度区域定义，比光源色度更宽松。
\end{enumerate}

\section{位置与几何可见度}

\subsection{装置安装数量}

\begin{table}[H]
\centering
\caption{各装置安装数量对比}
\begin{tabular}{|l|c|c|c|}
\hline
\textbf{装置类型} & \textbf{\versionAShort 数量} & \textbf{\versionBShort 数量} & \textbf{\versionCShort 数量} \\
\hline
远光灯          & 2 或 4 盏  & 同左                & 2 or 4 \\
\hline
近光灯          & 2 盏      & 同左                & 2 \\
\hline
前雾灯             & 2 盏      & 同左                & 2 \\
\hline
后雾灯             & 1 或 2 盏  & 同左                & 1 or 2 \\
\hline
倒车灯             & 1 或 2 盏  & 同左                & 1 or 2 \\
\hline
前转向灯           & 2 盏      & 同左                & 2 \\
\hline
后转向灯           & 2 盏      & 同左                & 2 \\
\hline
侧转向灯           & 2 或 4 盏  & 强化规定：增加最低数量 & 2 or more \\
\hline
危险警告           & 与转向灯共用 & 同左 & Same as indicators \\
\hline
制动灯             & 2 盏（S3 类 1 盏） & 同左，新增 ESS 单独规定 & 2 (S3 may be 1) \\
\hline
牌照灯             & 1 或多盏   & 同左                & 1 or more \\
\hline
前位置灯           & 2 盏      & 同左                & 2 \\
\hline
后位置灯           & 2 盏      & 同左                & 2 \\
\hline
DRL                & 2 盏      & 强化：2 盏必须       & 2 mandatory \\
\hline
侧标志灯           & 多选 & 较短车辆强制 & Optional/mandatory \\
\hline
前示廓灯           & 2 盏      & 同左                & 2 \\
\hline
后示廓灯           & 2 盏      & 同左                & 2 \\
\hline
后反射器           & 2 或 4 盏  & 同左                & 2 or 4 \\
\hline
侧反射器           & 多盏      & 强化规定：间距要求   & Multiple, with spacing \\
\hline
前反射器           & 2 盏（非强制） & 同左，新增"低速电动车强制" & Optional \\
\hline
ADB              & 高限       & 高限（强化）            & Allowed \\
\hline
\end{tabular}
\end{table}

\section{电气连接与控制}

\subsection{电路连接要求}

\begin{table}[H]
\centering
\caption{电路连接要求对比}
\small
\begin{tabularx}{\textwidth}{|l|X|X|X|}
\hline
\textbf{项目} & \textbf{\versionAShort} & \textbf{\versionBShort} & \textbf{\versionCShort} \\
\hline
前照灯电路       & 独立控制，一般开关或自动   & 强化，新增"ADB 单独控制"要求 & Independent control \\
\hline
近光/远光切换    & 一般规定                & 强化，新增"ADB 切换逻辑"     & Switching logic \\
\hline
示廓灯电路       & 前后独立                & 强化                     & Independent \\
\hline
制动灯电路       & 与制动踏板联动            & 强化，新增"ESS 独立电路"    & Linked with brake \\
\hline
转向灯电路       & 与转向信号联动            & 强化，新增"故障指示"要求   & Linked with indicator \\
\hline
\end{tabularx}
\end{table}

\subsection{故障指示要求}

\begin{table}[H]
\centering
\caption{故障指示要求对比}
\begin{tabular}{|l|c|c|c|}
\hline
\textbf{项目} & \textbf{\versionAShort} & \textbf{\versionBShort} & \textbf{\versionCShort} \\
\hline
转向灯故障指示       & 一般要求 & 强化，新增"对 LED 的故障定义" & 详细规定 \\
\hline
制动灯故障指示       & 一般要求 & 强化，新增"对 ESS 的故障定义"  & 详细规定 \\
\hline
前照灯故障指示       & 一般要求 & 强化，新增"ADB 故障降级模式" & 详细规定 \\
\hline
DRL 故障指示        & 一般要求 & 强化                          & 详细规定 \\
\hline
后雾灯故障指示      & 一般要求 & 强化                          & 详细规定 \\
\hline
\end{tabular}
\end{table}

\subsection{互连系统}

\begin{table}[H]
\centering
\caption{互连系统要求对比}
\small
\begin{tabularx}{\textwidth}{|l|X|X|X|}
\hline
\textbf{项目} & \textbf{\versionAShort} & \textbf{\versionBShort} & \textbf{\versionCShort} \\
\hline
挂车/牵引车连接   & 基本规定            & 强化，新增"对 LED 灯光信号设备的联调"   & 详细规定 \\
\hline
多装置联动        & 一般规定            & 强化，新增"互不干扰的逻辑要求"   & 详细规定 \\
\hline
网络通信          & 未明确             & 新增 CAN/LIN 总线通信要求 & 详细规定 \\
\hline
远程诊断          & 未明确             & 新增 远程诊断接口要求   & 详细规定 \\
\hline
\end{tabularx}
\end{table}

\section{特殊要求对比}

\subsection{恶劣环境适应性}

\begin{table}[H]
\centering
\caption{恶劣环境适应性要求对比}
\begin{tabular}{|l|c|c|c|}
\hline
\textbf{项目} & \textbf{\versionAShort} & \textbf{\versionBShort} & \textbf{\versionCShort} \\
\hline
温度范围         & -40°C 至 +85°C & 同左，强化上限值 & -40°C to +85°C \\
\hline
湿度             & 95\% RH         & 同左              & 95\% RH \\
\hline
振动             & 一般规定        & 强化，详细规定     & 详细规定 \\
\hline
防水防尘         & IPX4/IP5X       & 强化，新增"ADB 模块"IP 等级 & IPX4/IP5X \\
\hline
盐雾             & 一般规定        & 强化              & 详细规定 \\
\hline
\end{tabular}
\end{table}

\subsection{EMC 要求（草案版强化）}

\begin{table}[H]
\centering
\caption{EMC 要求对比}
\small
\begin{tabularx}{\textwidth}{|l|X|X|X|}
\hline
\textbf{项目} & \textbf{\versionAShort} & \textbf{\versionBShort} & \textbf{\versionCShort} \\
\hline
辐射发射限制       & 一般引用其他标准    & 强化，新增"对 ADB 的具体限值"  & 详细规定 \\
\hline
传导发射限制       & 一般引用其他标准    & 强化，具体限值明确 & 详细规定 \\
\hline
抗扰度要求         & 一般引用其他标准    & 强化                  & 详细规定 \\
\hline
静电放电           & 一般引用其他标准    & 强化                  & 详细规定 \\
\hline
瞬态传导           & 一般引用其他标准    & 强化                  & 详细规定 \\
\hline
\end{tabularx}
\end{table}

\subsection{光学性能特殊要求}

\begin{table}[H]
\centering
\caption{光学性能特殊要求对比}
\begin{tabular}{|l|c|c|c|}
\hline
\textbf{项目} & \textbf{\versionAShort} & \textbf{\versionBShort} & \textbf{\versionCShort} \\
\hline
近光防眩目       & 明暗截止线要求     & 强化，新增"ADB 防眩目模式" & 详细规定 \\
\hline
远光均匀性       & 一般规定          & 强化                  & 详细规定 \\
\hline
转向灯闪烁       & 60--120 次/分 & 同左，新增"对不同类型的具体规定" & 详细规定 \\
\hline
ESS 闪烁       & 4 Hz--6 Hz      & 强化                  & 4 Hz--6 Hz \\
\hline
\end{tabular}
\end{table}

\section{小结}

\begin{enumerate}[leftmargin=2em]
    \item 草案版在一般要求上向 UN R48-09 进一步靠拢，特别是在 EMC、智能控制、互连系统等方面；
    \item 草案版保留了中国对部分项目的特色规定，体现了国情的考量；
    \item 草案版对 LED 灯具的热管理、降级模式等作出了新增规定；
    \item 草案版的色度坐标与现行版、UN R48 基本保持一致，确保了互换性。
\end{enumerate}